\documentclass{article}
\usepackage[T1]{fontenc}
\usepackage{lmodern}
\usepackage{hyperref}

\title{The \texttt{pdfannex-fixture} skeleton adapter}
\author{pdfannex-adapter-template contributors}
\date{Version 0.1.0} % # x-release-please-version

\begin{document}
\maketitle

\section{Purpose}

\texttt{pdfannex-fixture} adds a command for documents served by a plain
local directory tree. It is distributed separately from the core
\texttt{pdfannex} package. The core package does not access the network or
know about this adapter.

This is a minimal skeleton adapter shipped with
\href{https://github.com/nocona71/pdfannex-adapter-template}
{\texttt{pdfannex-adapter-template}} so that the template is runnable and
testable end to end. \textbf{It is not a production service adapter.} It has
no network access, authentication or fault injection. Start a new adapter by
copying the template and replacing the resolver's service-specific
functions with calls to your real service's API.

\section{Requirements}

Use \texttt{pdfannex} dated 2026/05/01 or later and \texttt{texlua}, provided
by TeX Live. The resolver runs only during explicit CLI operations, never
while compiling.

\section{Use}

\begin{verbatim}
\usepackage{pdfannex-fixture}
...
\includedoc[label=memo]{memo}{Project memo}
\end{verbatim}

Set \texttt{PDFANNEX\_FIXTURE\_ROOT} to a directory containing
\texttt{<id>/<rev>.pdf} files. Without \texttt{rev}, the resolver pins the
latest revision during \texttt{pdfannex prepare}.

\section{See also}

The \href{https://github.com/nocona71/pdfannex-adapter-template}{template
repository} documents the full Resolver Protocol 1, the conformance check,
and the test/packaging harness this adapter exercises.

\end{document}
